%=====================================================================
\chapter{Answers to the Checkpoints}\label{app:answers}
%=====================================================================
\normalfigures

Three answers per chapter, in order. Attempt each checkpoint from memory before
reading these. Several are judgement calls rather than facts, and say so; where
an answer quotes a number, the number comes from the chapter's harness or from
the seeded database, and you can reproduce it.

% The slide pipeline parses this appendix: \ansch{<number> --- <title>} must be
% followed immediately by an ans environment of exactly three items.
\newcommand{\ansch}[1]{\subsection*{\color{primary}Chapter #1}}
\newenvironment{ans}
  {\begin{enumerate}[leftmargin=*, itemsep=4pt, topsep=3pt,
                     label=\textbf{\arabic*.}]}
  {\end{enumerate}}

%---------------------------------------------------------------------
\ansch{1 --- Why Not Just Files?}
\begin{ans}
  \item \textbf{It escapes most of them, for now.} With one clerk per file and
        no shared data there is no redundancy, no inconsistency between copies
        and no concurrent access. It still has \emph{isolation} (``which
        borrowers are also in arrears?'' needs both files), \emph{integrity}
        (rules live in each clerk's head) and \emph{atomicity} (a half-saved
        file). The answer changes the moment any fact --- a student's name ---
        is recorded in both files, or a second clerk edits either one.
  \item A \textbf{database} is the data together with its description; a
        \textbf{DBMS} is the software that manages it. ``We bought a database''
        might mean a product (SQLite, Oracle) with no data in it, or a
        populated set of tables, or a whole system --- three purchases with
        different costs, owners and risks. You cannot install, back up or
        migrate it until you know which.
  \item Disk is cheap; \textbf{disagreement} is not. Redundancy's real cost is
        \emph{inconsistency}: two copies of one fact that can diverge. Example:
        a student's name changes and is corrected in \texttt{students.csv} but
        not \texttt{fees.csv}. Each file is internally correct; the system now
        holds two names for one person, and the fee receipt goes to someone
        who no longer exists.
\end{ans}

\ansch{2 --- Three Levels, and Why They Are Separate}
\begin{ans}
  \item Adding an index: \textbf{internal}, physical independence. Adding a
        column: \textbf{conceptual}, logical independence (programs that never
        named it carry on). Renaming a column: \textbf{conceptual}, and logical
        independence \emph{fails} unless a view keeps the old name. A faster
        disk: \textbf{internal}, physical independence. A view with no home
        addresses: \textbf{external}; no independence is involved --- it is
        the external level being used.
  \item A schema is the \emph{form}; an instance is what is \emph{written on
        it}. A blank admission form is the schema: it says there will be a
        name, a date of birth, a programme. This year's 4{,}000 filled-in
        forms are the instance. Reprinting the form with a new box is a schema
        change and is rare; filling one in is an instance change and happens
        every minute.
  \item The \textbf{external level} is affected: views can define it, but
        nothing can confine a user to one, because SQLite has no users. You can
        still write the views and have programs use only them. You cannot stop
        anyone who can open the file from reading the base tables. In an
        application the program itself becomes the gatekeeper --- it holds the
        file, authenticates people, and only ever queries the views --- which
        is enforcement by convention, not by the database.
\end{ans}

\ansch{3 --- What Is Inside a DBMS, and Who Does What}
\begin{ans}
  \item The \textbf{optimiser} (a different plan --- perhaps statistics were
        refreshed, or an index dropped) and the \textbf{buffer manager} (Monday
        the pages were in memory; Friday they were not). Tell them apart with
        \sqlkw{EXPLAIN QUERY PLAN}: if the plan changed, it is the optimiser;
        if the plan is identical, run the query twice --- a fast second run
        means the first was paying for disk reads.
  \item \textbf{Advantage:} any tool can discover a database it has never
        seen with an ordinary \sqlkw{SELECT}, and the schema is backed up with
        the data. \textbf{Problem:} the dictionary can be damaged by the same
        failures as the data, and in SQLite anyone who can write the file can
        in principle corrupt \texttt{sqlite\_master} and with it every table.
  \item DBA: when to back up. Designer: what the keys are. Programmer: how a
        query is parameterised. End user: what question is asked. In this
        course the designer, programmer and DBA decisions are all yours; the
        DBA decisions that need users and privileges simply cannot be made,
        and backup becomes \texttt{.backup} run by you.
\end{ans}

\ansch{4 --- Your First Database, in One Command}
\begin{ans}
  \item Page 1 (4{,}096 bytes) holds the 100-byte database header and the
        schema table \texttt{sqlite\_master} --- the data dictionary. Page 2
        (4{,}096 bytes) is the root page of the new table's B-tree, empty.
        Fewer is impossible: storage is allocated in whole pages, the
        dictionary must exist for the file to describe itself, and the table
        must have a root.
  \item From 1{,}000 to 10{,}000 rows the file grew $262{,}144 - 32{,}768 =
        229{,}376$ bytes, about 25.5 bytes a row, so 2{,}000 rows predicts
        $32{,}768 + 25{,}500 \approx 58{,}300$, rounded up to a whole page:
        61{,}440 (15 pages). \textbf{Measured: 57{,}344 bytes, 14 pages.} The
        naive guess --- double 32{,}768 --- gives 65{,}536 and is further off,
        because part of the file is fixed overhead (the header page and the
        tree's interior page) that does not grow with the rows.
  \item
\begin{lstlisting}[language=sql]
CREATE TABLE instructor (
  emp_id      TEXT PRIMARY KEY NOT NULL,
  name        TEXT NOT NULL,
  department  TEXT NOT NULL,
  designation TEXT
) STRICT;
\end{lstlisting}
        \texttt{emp\_id} is \textbf{text} because \texttt{E-0114} has
        a letter and a leading zero and is never used in arithmetic; it is
        \textbf{not null} because an identifier that is missing identifies
        nothing. \texttt{designation} may be null: a new appointment may not
        have one yet.
\end{ans}

\ansch{5 --- The Entity-Relationship Model}
\begin{ans}
  \item ``Almost every'' means some have two, and one column cannot hold two
        values that can be searched. The first student with two numbers forces
        either a lost number or a comma-separated list, which no query can
        search correctly. Instead, a \texttt{student\_phone(roll\_no, phone)}
        table with both columns in the key; it costs one join whenever a
        number is shown.
  \item Candidate keys: $\{$emp\_id$\}$ (issued uniquely by the university),
        $\{$cnic$\}$ (national identity, unique by law), $\{$email$\}$ (issued
        one per employee). \textbf{Primary key: emp\_id} --- unique, always
        present (CNIC may be missing for foreign staff), never changing (email
        changes with a name or a reorganisation), short. \texttt{office\_no},
        \texttt{name} and \texttt{department} are not keys: offices are
        shared, names repeat, departments hold many people.
  \item A grade is a fact about \emph{a student in a particular section} ---
        neither alone determines it. On \ent{student} there is room for one
        grade per student for their whole degree. On \ent{section} there is
        room for one grade per section, the same for every student in it. Only
        the pairing of the two has a grade, and the pairing is the
        relationship.
\end{ans}

\ansch{6 --- Cardinality, Participation and Weak Entities}
\begin{ans}
  \item \textbf{It will refuse a section with no instructor.} A section is
        created in June for the autumn timetable; the lecturer is appointed in
        August. For two months the database will not accept the section ---
        and somebody will invent an instructor called ``TBA''.
  \item \textbf{M:N}, with \textbf{partial} participation of \ent{student}
        (a student may be enrolled in nothing) and \textbf{partial}
        participation of \ent{section} (a section may have no students). With
        $(1,N)$ at the section end, participation of \ent{section} becomes
        \textbf{total}: a section with no students could not exist, so a new
        section could not be created until its first enrolment arrived.
  \item Section numbers repeat across courses --- every course has a section
        1 --- so \texttt{section\_no} identifies a section only within its
        owner. Partial key: \texttt{section\_no}. Full key: $\{$course\_code,
        semester, section\_no$\}$. With a surrogate \texttt{section\_id} you
        must still declare \sqlkw{UNIQUE} (course\_code, semester,
        section\_no); forget it and the same section can be created twice,
        with students split between the copies.
\end{ans}

\ansch{7 --- The Relational Model}
\begin{ans}
  \item \textbf{Degree 5} counts attributes and belongs to the schema;
        \textbf{cardinality 200} counts tuples and belongs to the instance.
        Cardinality changes without a schema change, with every
        \sqlkw{INSERT} or \sqlkw{DELETE}. Degree changes only with a schema
        change: \sqlkw{ALTER TABLE} \dots\ \sqlkw{ADD} or \sqlkw{DROP COLUMN}.
  \item \texttt{SELECT COUNT(*) FROM enrolment WHERE roll\_no = ?} --- if the
        same enrolment can be inserted twice, a student's course count is
        inflated. Declare the composite \sqlkw{PRIMARY KEY} (roll\_no,
        course\_code, semester, section\_no), which makes the duplicate
        impossible.
  \item One null marker carries both meanings, and SQLite has no second kind.
        \textbf{Design 1:} a \texttt{grade\_status} column (\texttt{'not sat'},
        \texttt{'lost'}, \texttt{'graded'}) beside a nullable grade --- one
        column more, and a rule that the two agree. \textbf{Design 2:} a
        separate \texttt{exam\_incident} table recording lost papers --- no
        nulls overloaded, but a join to tell the two cases apart.
\end{ans}

\ansch{8 --- Keys and Integrity Constraints}
\begin{ans}
  \item \texttt{PRAGMA foreign\_keys;} tells you whether \emph{this
        connection} enforces --- not whether anyone else's does, or ever did.
        \texttt{PRAGMA foreign\_key\_list(t);} shows the declarations exist.
        \texttt{PRAGMA foreign\_key\_check;} lists every orphan already
        present: rows means enforcement was off at some point; none means no
        violation survives, \emph{not} that enforcement was on. Nothing in the
        file records whether it ever was.
  \item \ent{section}$\rightarrow$\ent{course}: \sqlkw{CASCADE} --- a
        section is a weak entity and means nothing without its course.
        \ent{section}$\rightarrow$\ent{instructor}: \sqlkw{SET NULL} --- the
        section survives a resignation; participation is partial.
        \ent{enrolment}$\rightarrow$\ent{student}: \sqlkw{RESTRICT} --- a
        student with a grade history must never be silently deleted.
  \item It will accept a \textbf{null} primary key and store it, violating
        \textbf{entity integrity}. Fix it with \sqlkw{NOT NULL} on the column,
        or \sqlkw{STRICT} on the table --- measured in \chref{constraints},
        each one alone refuses the null.
\end{ans}

\ansch{9 --- Mapping an ER Diagram to Tables}
\begin{ans}
  \item On \ent{department}, as \texttt{head\_emp\_id}: every department must
        have a head (total), but most instructors head nothing (partial), so
        the totally-participating side can declare the column \sqlkw{NOT NULL}.
        On \ent{instructor} instead it would be \texttt{heads\_dept\_code},
        null for all but eight of a hundred instructors.
  \item \texttt{roll\_no}: who. \texttt{course\_code}: which course.
        \texttt{semester}: when. \texttt{section\_no}: which group within that
        offering. Without \texttt{semester} a student could not appear in the
        same course and section number in two terms --- so repeating a failed
        course becomes impossible, or overwrites the old grade.
  \item The table \textit{prerequisite}, whose primary key is the pair and
        both of whose columns reference \textit{course}; the unstated rule is
        the \sqlkw{CHECK}:
\begin{lstlisting}[language=sql]
CREATE TABLE prerequisite (
  course_code TEXT NOT NULL REFERENCES course(course_code),
  requires    TEXT NOT NULL REFERENCES course(course_code),
  PRIMARY KEY (course_code, requires),
  CHECK (course_code <> requires)    -- nothing requires itself
) STRICT;
\end{lstlisting}
 (A two-step cycle needs a trigger;
        \chref{subqueries} shows why it also needs \sqlkw{UNION}.)
\end{ans}

\ansch{10 --- SQL: Creating the Schema}
\begin{ans}
  \item The marks were 95, 72, \texttt{'absent'}, 88 and 100. Affinity left
        \texttt{'absent'} as text; \sqlkw{AVG} read it as 0, so the sum was
        $95 + 72 + 0 + 88 + 100 = 355$ and, because it is not null, the count
        was 5: $355 / 5 = 71$. Proof: \texttt{SELECT typeof(score), COUNT(*)
        FROM marks GROUP BY 1;} shows the text row, and \texttt{SELECT
        SUM(score), COUNT(score) FROM marks;} returns 355 and 5.
  \item A column constraint follows one column and can name only it; a table
        constraint follows all the columns and can name several. In
        \texttt{v1\_create.sql}: the composite \sqlkw{PRIMARY KEY} of
        \textit{section} (three columns), and the three-column
        \sqlkw{FOREIGN KEY} of \textit{enrolment} --- neither belongs to one
        column, so neither can be written beside one.
  \item \textbf{Compatibility}: other tools and older SQLite versions (before
        3.37) cannot read a \sqlkw{STRICT} table, and it forbids type names
        like \texttt{VARCHAR} or \texttt{DATETIME} that ported SQL uses. That
        is a real cost when the file is shared with old software; for this
        course's database it does not apply, and the measured cost was zero.
\end{ans}

\ansch{11 --- Changing a Schema That Is Already Full}
\begin{ans}
  \item Nullable \texttt{email}: \textbf{cheap} --- no existing row changes.
        \texttt{email} \sqlkw{NOT NULL}: \textbf{expensive} --- every row needs
        a value, so backfill and rebuild. \texttt{semester} text to integer:
        \textbf{expensive} --- a rebuild, and step~0 must find any value that
        will not convert. Splitting a whole address: \textbf{impossible} to do
        exactly --- the parts were never recorded separately, and parsing free
        text guesses.
  \item After the migration, backfilled values are indistinguishable from
        values entered since; a rollback that drops the column cannot tell
        which were invented, and one that keeps them keeps invented data. To
        make it reversible, record at migration time which rows were
        backfilled --- a \texttt{migration\_backfill(roll\_no)} table, or a
        flag column --- so the reverse can remove exactly those.
  \item If the migration fails after \sqlkw{COMMIT} but the version is set
        afterwards, that is merely late. The real danger is the reverse: set
        after a migration that failed and was rolled back, and the file claims
        to be version 2 while holding the version 1 schema. Every later
        migration then assumes columns that are not there.
\end{ans}

\ansch{12 --- Relational Algebra}
\begin{ans}
  \item A relation is a \emph{set}, so $\pi$ removes duplicates: four
        programmes, four tuples. A SQL result is a \emph{bag} and keeps every
        row. \sqlkw{DISTINCT} reconciles them:
        \texttt{SELECT DISTINCT programme FROM student} returns 4.
  \item $\pi_{roll}(\sigma_{semester=2}(\mathit{student})) -
        \pi_{roll}(\sigma_{semester='FALL2026'}(\mathit{enrolment}))$, that is:
\begin{lstlisting}[language=sql]
SELECT roll_no FROM student   WHERE semester = 2
EXCEPT
SELECT roll_no FROM enrolment WHERE semester = 'FALL2026';
\end{lstlisting}
        ``Nothing'' is the difference. The tempting wrong version, a join with
        \texttt{e.semester <> 'FALL2026'}, returns second-year students enrolled
        in \emph{some other} term --- including those also enrolled in
        FALL2026. Note the two meanings of \texttt{semester}: \chref{joins}
        returns to them.
  \item $4{,}000 \times 60{,}037 = 240{,}148{,}000$ intermediate tuples, of
        which 60{,}037 survive the selection. \textbf{Closure} --- every
        intermediate result is a relation --- is what makes algebraic laws such
        as ``select before you multiply'' valid, so the optimiser may compute
        the join by index lookups and still return exactly the same relation.
\end{ans}

\ansch{13 --- Joins: the Algebra and the SQL}
\begin{ans}
  \item The two tables share \textbf{two} column names: \texttt{roll\_no} and
        \texttt{semester}. \sqlkw{NATURAL JOIN} required both to be equal, and
        a year of study (\texttt{8}) never equals a term label
        (\texttt{'FALL2025'}). In the query: \texttt{JOIN \dots\ USING
        (roll\_no)} or \sqlkw{ON}. In the schema: rename
        \texttt{student.semester} to \texttt{year\_of\_study} --- it is the one
        whose name is misleading.
  \item They are the 73 sections whose \texttt{emp\_id} is null.
        \texttt{SELECT s.* FROM section s LEFT JOIN instructor i ON
        s.emp\_id = i.emp\_id WHERE i.emp\_id IS NULL;} The decision was making
        \ent{section}'s participation in \ent{teaches} \textbf{partial} in
        \chref{cardinality} --- deliberately, so a section can exist before its
        lecturer is appointed.
  \item Without an index, every pair of 4{,}000 students was compared ---
        about 8 million --- to find 28{,}517 matches; with one, only the
        matches were visited: 108$\times$ faster. Growth ratios mislead because
        the number of \emph{matches} itself grew four-fold per doubling (names
        come from a fixed pool), so even the indexed join looked quadratic down
        its column. Only a comparison at equal size isolates search cost.
\end{ans}

\ansch{14 --- Selecting Rows, and the Trouble with NULL}
\begin{ans}
  \item $60{,}037 - 7{,}268 - 32{,}777 = 19{,}992$ rows: the ungraded FALL2026
        enrolments, whose grade is null. For them both \texttt{= 'A'} and
        \texttt{<> 'A'} are \unk{}, and \sqlkw{WHERE} keeps only \emph{true}.
        \texttt{WHERE grade <> 'A' OR grade IS NULL} returns all 52{,}769 that
        are not an A (in SQLite \texttt{WHERE grade IS NOT 'A'} says the same).
  \item \texttt{x NOT IN (a, NULL)} means \texttt{x <> a AND x <> NULL}. The
        second comparison is \unk{}. If the first is false the whole is false;
        if the first is true, \texttt{T AND U = U}. So for every $x$ the
        condition is false or \unk{} --- never true --- and \sqlkw{WHERE}
        discards every row.
  \item It decides that \textbf{a missing mark is a zero} --- an absent
        student failed. Scores 80, 90, null: \sqlkw{AVG} gives
        $(80+90)/2 = 85$; \sqlkw{AVG}(\sqlkw{COALESCE}) gives $(80+90+0)/3 =
        56.7$. Write \sqlkw{AVG}(score), and report the count of missing marks
        beside it, so the decision about them is made by the examinations
        office rather than by a function.
\end{ans}

% Plain text: the slide parser reads this heading with [^}]*, so a
% brace inside the title would end it early and lose all three answers.
\ansch{15 --- Sub-queries, EXISTS and Set Operations}
\begin{ans}
  \item The \sqlkw{IN} sub-query of step~2 is \textbf{uncorrelated}. Its
        \sqlkw{EXISTS} is \textbf{correlated}, on \texttt{s.roll\_no}. In
        step~3, \sqlkw{NOT IN} and the two \sqlkw{EXCEPT} branches are
        uncorrelated; \sqlkw{NOT EXISTS} is correlated, on \texttt{i.emp\_id}.
        Steps~4 and~5 are correlated, on \texttt{c.course\_code}. The recursive
        CTE's step refers to itself, not to an outer query.
  \item \sqlkw{NOT IN} returned 0. Its list contained the 73 null
        \texttt{emp\_id}s, so each instructor's test included \texttt{x <>
        NULL}, which is \unk{} --- and no row's condition could be true.
        \sqlkw{EXCEPT} is a set operation, not a comparison: it removes a
        left-hand row only if an equal row appears on the right, and a null on
        the right could only remove a null on the left. \texttt{instructor.emp\_id}
        has none.
  \item Without an index, each of 48 courses scanned all 60{,}037
        enrolments --- 2.9 million row visits --- while \sqlkw{GROUP BY} read
        them once. With an index, each course became one index search and the
        correlated form read only what it counted, while \sqlkw{GROUP BY} still
        read everything. The advice is really ``avoid \emph{unindexed}
        correlated sub-queries''; correlation itself costs nothing.
\end{ans}

\ansch{16 --- Grouping, Aggregation and Window Functions}
\begin{ans}
  \item 2.383 divides the points by the \textbf{40{,}045} graded enrolments;
        1.590 divides the same points by all \textbf{60{,}037}, counting each
        ungraded enrolment as a zero. Publish 2.383, labelled as the average
        over graded courses, with the 19{,}992 in progress reported separately:
        an unfinished course is not an F.
  \item \texttt{course\_code LIKE 'CS-\%'} mentions only a grouping column, so
        it gives the same answer whether rows are removed before grouping or
        whole groups after --- and SQLite moved it into \sqlkw{WHERE} itself.
        \texttt{COUNT(*) > 1000} needs the count, which does not exist until
        grouping is finished, so it can only be tested after. Write
        \sqlkw{WHERE} because it states what you mean; the rewrite is an
        optimisation, not a promise.
  \item ``Each student's rank within their programme.''
        \texttt{SELECT roll\_no, programme, gpa, RANK() OVER (PARTITION BY
        programme ORDER BY gpa DESC) FROM per\_student;} \sqlkw{GROUP BY}
        programme would collapse each programme to one row and lose every
        student; the window keeps all of them and adds the rank.
\end{ans}

\ansch{17 --- Functional Dependencies}
\begin{ans}
  \item $\{$emp\_id$\}^{+}$: start \{emp\_id\}; \fd{emp\_id}{instructor} adds
        instructor; nothing else applies: \textbf{\{emp\_id, instructor\}}.
        $\{$course\_code, semester, section\_no$\}^{+}$: pass 1 adds room and
        emp\_id (the section rule) and title, credit\_hours, dept\_code (the
        course rule); pass 2 adds instructor (via emp\_id) and dept\_name (via
        dept\_code); pass 3 adds nothing. Ten attributes; not a superkey ---
        roll\_no, name, programme and grade are missing.
  \item $\{$roll\_no, course\_code$\}^{+}$: pass 1 adds name, programme (from
        roll\_no) and title, credit\_hours, dept\_code (from course\_code);
        pass 2 adds dept\_name. \textbf{dept\_name is in the closure, so the
        dependency follows.}
  \item Step 1: already single right sides. Step 2: in \fd{AC}{D}, is $C$
        extraneous? $\{A\}^{+} = \{A, B, C\}$ by \fd{A}{B} and \fd{A}{C}, then
        \fd{AC}{D} gives $D$ --- so \fd{A}{D} holds and $C$ is extraneous.
        Step 3: \fd{A}{C} is redundant (\fd{A}{B}, \fd{B}{C}). Result
        $\{$\fd{A}{B}, \fd{B}{C}, \fd{A}{D}$\}$; $\{A\}^{+} = \{A,B,C,D\}$ in
        both.
\end{ans}

\ansch{18 --- Normal Forms: 1NF to BCNF}
\begin{ans}
  \item \textbf{Update}: renaming department CS changed 7{,}869 rows (one in
        \textit{department}), and a half-done rename left two names --- caused
        by \fd{dept\_code}{dept\_name}. \textbf{Insertion}: a new course with no
        enrolments had no row to live in --- \fd{course\_code}{title, \dots}
        with a key that needs a student. \textbf{Deletion}: removing PH-350's
        1{,}122 enrolments removed the course --- the same dependency.
  \item Key $\{A, B\}$. \fd{A}{C} is \textbf{partial} (violates 2NF);
        \fd{C}{D} is \textbf{transitive} (violates 3NF). So R is in 1NF only.
        Deepest first: on \fd{C}{D}, $R_1(\underline{C}, D)$, remainder
        $(A, B, C)$; on \fd{A}{C}, $R_2(\underline{A}, C)$, remainder
        $R_3(\underline{A}, \underline{B})$. All three are in BCNF.
  \item The transcript was a \emph{lookup}: the normalised \textit{enrolment}
        could search its primary key by roll\_no, while \textit{report} had no
        index at all and was read end to end. The department report was a
        \emph{full scan}, cheaper over one table than over three. With an index
        on roll\_no the transcript result changes to \textbf{a tie} (0.08 against
        0.09\,ms); the full-scan result does not change.
\end{ans}

\ansch{19 --- Decomposition, and Denormalising on Purpose}
\begin{ans}
  \item $R_1 \cap R_2 = \{$course\_code$\}$. Its closure (under the rules
        these two tables can express) does not contain roll\_no or grade, so
        it is a superkey of neither: \textbf{lossy}. Joining on course\_code
        pairs every student who took a course with every grade given in it ---
        182{,}920 rows from 39{,}152 true facts.
  \item Split \textit{course} into \texttt{course\_dept(course\_code,
        dept\_code)} and \texttt{course\_dname(course\_code, dept\_name)}:
        lossless (course\_code is the key of both), but \fd{dept\_code}{dept\_name}
        lies in neither. Inserting \texttt{('CS-150', 'Computing')} when
        another CS course says \texttt{'Computer Science'} is accepted ---
        measured in \chref{decomposition}.
  \item Pays off when reads per write exceed $3.679 / 69.73 = 0.053$: one
        read in every 19 writes. On a 1\,ms-commit SSD, writes get cheaper in
        general, and the part of the 3.679\,ms that came from making a second
        table's page durable at each commit shrinks; \chref{views} measured the
        trigger's own work at about a hundredth of a millisecond. The added
        cost per write falls, so the
        break-even falls too --- the summary table pays off at even fewer
        reads.
\end{ans}

\ansch{20 --- Views, Triggers and Keeping Data Honest}
\begin{ans}
  \item SQLite stored only the view's \sqlkw{SELECT} text, substituted it into
        the query that named it, and compiled the result. What runs is the
        hand-written query with renumbered cursors. Nothing was stored, so
        nothing can be out of date: each use is the query run now.
  \item Choose the table when measured read cost matters --- here 700$\times$
        --- and reads greatly outnumber writes. Alongside it: triggers on
        \sqlkw{INSERT}, \sqlkw{DELETE} \emph{and} \sqlkw{UPDATE} of every
        table it summarises; a query comparing it with the equivalent view,
        run on a schedule; and a script that rebuilds it from the base tables
        when they disagree.
  \item
\begin{lstlisting}[language=sql]
CREATE TRIGGER max_six_upd
BEFORE UPDATE OF semester, roll_no ON enrolment
WHEN (SELECT COUNT(*) FROM enrolment
       WHERE roll_no = NEW.roll_no
         AND semester = NEW.semester) >= 6
BEGIN SELECT RAISE(ABORT, 'at most six sections a term'); END;
\end{lstlisting}
        It uses \texttt{NEW},
        because the limit concerns the term the enrolment is moving
        \emph{into}; \texttt{OLD} is the term it is leaving.
\end{ans}

\ansch{21 --- Indexes, and Where They Stop Helping}
\begin{ans}
  \item Each page of a B-tree divides its range among a few hundred children,
        so three levels reach millions of entries and a lookup reads one page
        per level --- three or four reads at 10{,}000 rows or a million.
        Hence 0.053 against 0.057\,ms. A scan reads every page, and a hundred
        times the rows is roughly a hundred times the pages: 1.37 to
        199.90\,ms.
  \item It lacked the \textbf{distribution} of \textit{k}: \texttt{sqlite\_stat1}
        recorded only the row count and one row per key, so every range looked
        selective. That would come from \textbf{STAT4} histograms, which the
        SQLite shipped with Python does not include. To force the scan: write
        \texttt{+k} in the condition, or \texttt{FROM t NOT INDEXED}.
  \item \texttt{roll\_no = ?}: \textbf{yes} --- the leftmost column.
        \texttt{course\_code = ?}: \textbf{no} --- course\_code values are
        scattered through every roll\_no. Both: \textbf{yes} --- the full
        prefix, the most selective use.
\end{ans}

\ansch{22 --- Transactions, Concurrency and Recovery}
\begin{ans}
  \item \textbf{Atomicity}: the failed move left a student in nothing under
        autocommit, unchanged inside a transaction. \textbf{Consistency}:
        constraints from \chref{constraints} refuse any end state that breaks
        them. \textbf{Isolation}: \sqlkw{IMMEDIATE} counted both clerks where
        no transaction lost one. \textbf{Durability}: \texttt{synchronous =
        FULL} waited 23\,ms per commit for the disk; without it a power cut can
        lose a reported commit.
  \item \sqlkw{DEFERRED}: both took \sqlkw{SHARED} to read; A's write took
        \sqlkw{RESERVED}; B's write wanted \sqlkw{RESERVED} and was refused;
        A's \sqlkw{COMMIT} wanted \sqlkw{EXCLUSIVE}, which needs B's
        \sqlkw{SHARED} gone --- deadlock until B rolled back.
        \sqlkw{IMMEDIATE}: A took \sqlkw{RESERVED} at \sqlkw{BEGIN}, so B
        could not begin; after A committed, B began, read 41 and wrote 42.
  \item Each commit waited about 23\,ms for the disk to confirm; 200 waits is
        4.6\,s, and one transaction waits once. \texttt{synchronous = OFF}
        removes the wait by abandoning the guarantee --- and was still
        9$\times$ slower than batching. Used for results entry, a power cut in
        the last week of term could lose grades the clerks saw saved, or
        corrupt the file.
\end{ans}

\ansch{23 --- Security, NoSQL, and the Semester Project}
\begin{ans}
  \item The \textbf{single quote}. Pasted, it closed the string the programmer
        had opened, so \texttt{OR '1'='1'} became SQL and was true for all
        4{,}000 rows. Bound as a parameter, the value was never parsed: the
        quote was just a character in a roll number that no student has, so 0
        rows.
  \item
\begin{lstlisting}[language=sql]
GRANT SELECT         ON enrolment TO exams_office;
GRANT UPDATE (grade) ON enrolment TO exams_office;
\end{lstlisting}
        In SQLite
        the same office, if it can open the file, can read and change every
        table, including fees and addresses --- there are no users to grant to.
  \item ``All of one student's grades'': the \textbf{document} --- one lookup
        returns the whole record. ``A grades in CS-101'': \textbf{relational}
        --- one indexed condition on one table, against unpacking all 4{,}000
        documents (9.8 against 78.4\,ms on the seeded data). The model should
        follow the questions the system is asked most.
\end{ans}
